\documentclass[11pt,a4paper]{article}
\usepackage[utf8]{inputenc}
\usepackage[T1]{fontenc}
\usepackage[slovene]{babel}
\usepackage[margin=2cm]{geometry}
\usepackage{enumitem}
\setlist{nosep, leftmargin=*}
\setlength{\parindent}{0pt}
\setlength{\parskip}{4pt}
\pagestyle{empty}

\begin{document}

{\large\textbf{Prijava teme seminarske naloge --
NPMP, kvantni računalniki (QC)}}\\
Avtor: Matej Mlakar \hfill 

\vspace{4pt}
\textbf{Naslov:} Vpliv ravni optimizacije prevajalnika Qiskit
na reverzibilna vezja na šumnem kvantnem procesorju IBM

\textbf{Raziskovalno vprašanje:} Kako vpliva raven optimizacije
prevajalnika Qiskit (0--3) na število dvokubitnih vrat
reverzibilnih vezij in kako to vpliva na
verjetnost pravilnega rezultata na šumnem procesorju IBM?

\textbf{Pričakovani rezultat (hipoteza):}
\begin{itemize}
  \item Prehod z ravni 0 na raven 1
  močno zmanjša število dvokubitnih vrat (za več
  kot 30\,\%) in opazno izboljša pravilnost.
  \item Razlika med ravnema 2 in 3 je majhna (manj kot 10\,\% 
  v številu dvokubitnih vrat in manj kot 3\,\% v pravilnosti).
\end{itemize}

\textbf{Metrike:}
\begin{itemize}
  \item strukturne (pred prevajanjem in po njem):
  število kubitov, število dvokubitnih vrat, globina vezja,
  skupno število vrat, čas prevajanja;
  \item verjetnost pravilnega rezultata $P_{pravilno}$ = delež ponovitev
  s pravilnim izhodom, povprečen čez vse vhode
  vezja. Pravilen izhod določimo iz klasične pravilnostne
  tabele.
\end{itemize}

\textbf{Načrt poskusov:}
\begin{enumerate}
  \item Implementacija reverzibilnih vezij različnih velikosti: Toffolijeva
  in Fredkinova vrata, polni seštevalnik ter 2-bitni
  Cuccarov seštevalnik. Vhod pripravimo
  z vrati X, ločenimi z \texttt{barrier}, da
  optimizacija ne poseže v pripravo vhoda.
  \item Prevajanje vsakega vezja na ravneh 0--3
  z več semeni prevajalnika (\texttt{seed\_transpiler}) in
  zapis strukturnih metrik. Rezultate čez semena povprečimo.
  \item Serija A: privzeto prevajanje. Serija B: vse
  ravni z enako fiksno postavitvijo kubitov (\texttt{initial\_layout}),
  da ločimo vpliv postavitve od vpliva
  optimizacije vrat.
  \item Zagon na idealnem simulatorju (preverjanje pravilnosti,
  pričakovano 100\,\%) in na simulatorju s šumnim
  modelom procesorja IBM (vse kombinacije).
  \item Zagon izbrane podmnožice vezij na pravem procesorju
  IBM (Open Plan). Zabeležimo oznako procesorja, število
  ponovitev in ID-je nalog.
  \item Analiza: primerjava ravni, povezava med številom
  dvokubitnih vrat in pravilnostjo, grafi in tabele.
\end{enumerate}

\textbf{Potrebna orodja:} Python 3, Qiskit (prevajalnik,
\texttt{generate\_preset\_pass\_manager}), Qiskit Aer (idealni in šumni
simulator), Qiskit IBM Runtime (dostop do procesorja IBM
in šumni modeli \texttt{FakeBackend}), matplotlib in pandas
(analiza in grafi), GitHub (repozitorij), \LaTeX{} (poročilo).

\textbf{Razdelitev dela:} Seminarsko nalogo izdelujem sam,
zato vse dele (vezja, poskuse, analizo in poročilo)
pripravim sam.

\end{document}
